\chapter{Dark Pools et transactions de bloc}\label{ch:mx:dark-pools-and-blocks}

Un fonds de pension veut acheter deux millions d'actions d'un titre qui s'échange à hauteur de cinq millions par jour. Affiché dans un carnet d'ordres visible, l'ordre ferait bouger le cours avant
qu'un dixième de celui-ci ne soit exécuté~; envoyé dans un dark pool, il attend un vendeur, et chaque plateforme qu'il touche apprend un peu de ce qu'il cherche. Les plateformes obscures existent pour permettre
aux gros ordres de se rencontrer sans s'annoncer, et elles ne fonctionnent que dans la mesure où elles gardent le secret. Ce chapitre construit un dark pool et une
plateforme de transactions de bloc à côté d'un carnet visible dans le simulateur de bourse, mesure qui se fait exécuter dans l'obscurité, et traite un gros ordre de quatre façons pour chiffrer ce qui fuit.

\section{Croisement au milieu et quantités minimales}

Un dark pool (livre~1, chapitre~9) n'affiche rien et croise généralement au milieu du meilleur cours acheteur et du meilleur cours vendeur du marché visible, de sorte que les deux côtés
économisent la moitié de la fourchette par rapport à une négociation dans le carnet visible~: entre 20.00 et 20.02, un cent par action pour chacun, soit 100~\$ sur 10\,000 actions. Le prix est emprunté~:
la plateforme a besoin des cotations visibles, et ses croisements ne sont aussi à jour que sa copie de celles-ci. \texttt{firm.exchsim} donne désormais à une plateforme obscure cette copie~: un
message de contrôle \emph{cotation de référence} (N) transmet le meilleur cours acheteur et le meilleur cours vendeur de la plateforme visible au moteur obscur après chaque changement, et les \omterm{def:mx:order-types-and-their-uses:peg}{ordres indexés} au milieu
se calent sur elle (\texttt{midpoint\_dark\_pool(reference=\dots)}).

\begin{definition}[Quantité minimale d'exécution]\label{def:mx:dark-pools-and-blocks:minqty}
Une \emph{quantité minimale d'exécution}\index{quantité minimale d'exécution} (minimum execution quantity) sur un ordre obscur au repos est la plus petite exécution qu'il accepte~: un ordre entrant qui
ne peut pas s'échanger avec lui pour au moins cette quantité est ignoré par lui. Sur un ordre entrant exécuter-ou-annuler, elle est la plus petite exécution totale qui vaut la peine d'être prise.
\end{definition}

Un achat indexé au milieu de 5\,000 actions au repos, avec un minimum de 1\,000, ignore une vente de 300 actions et s'échange avec une vente de 1\,500 actions. Le minimum est un filtre sur
les contreparties~: les petits ordres qui sondent la plateforme, et les petits ordres des traders rapides en général, ne peuvent pas l'atteindre. C'est aussi un filtre sur les exécutions~: le
gros ordre s'échange moins.

\begin{definition}[Réseau de croisement]\label{def:mx:dark-pools-and-blocks:crossing}
Un \emph{réseau de croisement}\index{réseau de croisement} (crossing network) apparie les ordres d'achat et de vente à un prix pris ailleurs (le milieu visible, un cours de clôture, un
indice de référence), sans \omterm{def:mx:fragmentation-and-routing:discovery}{formation des prix} qui lui soit propre, soit en continu, soit à des moments programmés.
\end{definition}

\section{Ordres conditionnels et plateformes de blocs}

\begin{definition}[Indication d'intérêt]\label{def:mx:dark-pools-and-blocks:ioi}
Une \emph{indication d'intérêt}\index{indication d'intérêt} (IOI) est un message non contraignant par lequel un trader peut vouloir négocier un titre, avec un sens
et parfois une taille ou un prix, envoyé à une contrepartie ou à une plateforme pour trouver l'autre côté d'une grosse transaction.
\end{definition}

\begin{definition}[Ordre conditionnel]\label{def:mx:dark-pools-and-blocks:conditional}
Un \emph{ordre conditionnel}\index{ordre conditionnel} (conditional order) est une \omterm{def:mx:dark-pools-and-blocks:ioi}{indication d'intérêt} au repos sur une plateforme de blocs~: lorsqu'il rencontre un conditionnel contraire d'au moins
la taille minimale des deux côtés, la plateforme invite les deux traders à \emph{confirmer}, c'est-à-dire à envoyer un ordre ferme, dans un court délai. Si les deux le font, le bloc
croise~; si l'un des deux ne le fait pas, rien ne s'échange.
\end{definition}

\begin{definition}[Taux de confirmation]\label{def:mx:dark-pools-and-blocks:firmup}
Le \emph{taux de confirmation}\index{taux de confirmation} (firm-up rate) d'un participant est la part des invitations qu'il reçoit auxquelles il répond par un ordre ferme. Les plateformes de blocs
le suivent et restreignent les participants dont le taux est faible, car un conditionnel qui ne se confirme pas a appris que l'autre côté existe sans rien payer.
\end{definition}

Les conditionnels permettent à un ordre de reposer sur plusieurs plateformes de blocs et algorithmes à la fois sans risquer une double exécution, puisque seule une confirmation l'engage. Leur
faiblesse tient à leur définition~: une invitation signale aux deux côtés qu'un gros contradicteur existe. Une contrepartie qui n'avait jamais eu l'intention de négocier obtient l'information
gratuitement, raison pour laquelle le \omterm{def:mx:dark-pools-and-blocks:firmup}{taux de confirmation} est surveillé (\cref{fig:mx:dark-pools-and-blocks:conditional}). \texttt{firm.darkpool.ConditionalBook} implémente
l'appariement, les invitations, les délais de confirmation et les taux.

\begin{omfigure}
\begin{tikzpicture}[font=\footnotesize,box/.style={draw,rounded corners,align=center,minimum height=0.8cm}]
\node[box,fill=omDef!12] (b) at (0,1.2) {fonds\\ achat 20\,000\\ min 1\,000};
\node[box,fill=omDef!12] (s) at (0,-1.2) {vendeur\\ vente 8\,000\\ min 2\,000};
\node[box,fill=omProp!15] (v) at (4,0) {plateforme de blocs\\ appariement $\Rightarrow$ invitation};
\node[box,fill=omThm!12] (x) at (8.3,1) {les deux confirment~:\\ 8\,000 croisent au milieu};
\node[box,fill=omExo!12] (l) at (8.3,-1) {un seul ne confirme pas~:\\ aucune transaction, mais il le sait};
\draw[->] (b) -- node[above,sloped] {conditionnel} (v);
\draw[->] (s) -- node[below,sloped] {conditionnel} (v);
\draw[->] (v) -- (x);
\draw[->] (v) -- (l);
\end{tikzpicture}

\omcaption{Un bloc conditionnel. Deux ordres non contraignants se rencontrent, la plateforme invite les deux à confirmer dans un délai, et seuls deux ordres fermes s'échangent. Un
participant qui laisse expirer ses invitations apprend l'existence de l'autre côté sans négocier.}\label{fig:mx:dark-pools-and-blocks:conditional}
\end{omfigure}

\omcode{code/firm/darkpool/firm_darkpool.py}{70}{89}{Appariement des conditionnels par priorité temporelle et invitation des deux côtés~; un ordre n'appartient qu'à une invitation à la fois.}\label{lst:mx:dark-pools-and-blocks:invite}

\section{Qui négocie dans l'obscurité, et les preuves}

Zhu (2014) a expliqué qui passe dans l'obscurité. Les \omterm{def:mx:why-there-is-a-spread:informed}{traders informés} négocient dans le même sens les uns que les autres, se regroupent du côté chargé du marché et courent donc un
risque plus élevé de ne pas être exécutés dans l'obscurité, où l'autre côté doit venir à eux~; le marché visible leur convient mieux, le dark pool convient aux
traders non informés, et un dark pool peut concentrer l'information sur le marché visible. Comerton-Forde et Putnins (2015) ont constaté sur des données australiennes que les transactions obscures
sont moins informées que les transactions visibles, que de faibles niveaux de négociation obscure hors blocs sont bénins, voire bénéfiques pour l'efficience informationnelle, et que des niveaux élevés
lui sont nuisibles, et qu'aucune preuve ne montre que les transactions de bloc dans l'obscurité entravent la \omterm{def:mx:fragmentation-and-routing:discovery}{formation des prix}.

Le prix efficient $P^\ast$ du marché du chapitre évolue par des sauts exogènes et par l'impact permanent des \omterm{def:mx:the-limit-order-book:orders}{ordres au marché} non informés sur le
marché visible, $\lambda=0.0005$ tick par action (les ordres informés révèlent $P^\ast$, ils ne le font pas bouger). Les fournisseurs de liquidité cotent sur le marché visible~; les preneurs de bruit et
les preneurs informés essaient d'abord le dark pool la moitié du temps et envoient ce qui reste sur le marché visible~; les traders non informés patients laissent reposer des \omterm{def:mx:order-types-and-their-uses:peg}{ordres indexés} au milieu dans l'obscurité. Sur quatre
graines de 45 minutes, les preneurs informés qui ont essayé l'obscurité ont été exécutés à hauteur de 6.3~\% de leurs actions, les non informés à 17.5~\%~: les informés achètent quand
tout le monde informé achète, et les vendeurs de l'obscurité s'épuisent. Rester au repos dans l'obscurité abandonne la fourchette, mais ses exécutions ont été suivies d'un mouvement de $+0.11$ cent en
faveur du trader au repos sur 30 secondes, contre $-0.13$ cent pour les fournisseurs visibles, qui captent 0.51 cent de fourchette et subissent la sélection adverse sur celle-ci.

\section{Fuite d'information}

\begin{definition}[Fuite d'information]\label{def:mx:dark-pools-and-blocks:leakage}
La \emph{fuite d'information}\index{fuite d'information} (information leakage) est la part du coût d'un gros ordre qui provient de ce que d'autres traders en apprennent l'existence avant qu'il ne soit
terminé, à partir de ses exécutions, de ses cotations, de ses conditionnels ou de son routage, et de la négociation qui le précède.
\end{definition}

\begin{definition}[Ordre sonde]\label{def:mx:dark-pools-and-blocks:ping}
Un \emph{ordre sonde}\index{ordre sonde} (pinging order) est un petit \omterm{def:mx:order-types-and-their-uses:tif}{ordre exécuter-ou-annuler} envoyé à une plateforme obscure pour détecter la liquidité cachée~: une exécution révèle un
ordre contraire au repos à ce prix, et des exécutions répétées révèlent un ordre important.
\end{definition}

Un fonds achète 20\,000 actions à partir de la dixième minute de séance, sur 30 minutes, de quatre façons~: \emph{visible}, un \omterm{def:mx:the-limit-order-book:orders}{ordre au marché} de 200 toutes les 18 secondes~;
\emph{obscur}, un \omterm{def:mx:order-types-and-their-uses:peg}{ordre indexé} au milieu au repos pour l'ensemble de l'ordre~; \emph{obscur avec minimum}, le même \omterm{def:mx:order-types-and-their-uses:peg}{ordre indexé} avec un minimum de 1\,000~; \emph{bloc}, des conditionnels
dans une plateforme de blocs où des vendeurs naturels publient des conditionnels de 5\,000 à 15\,000 actions toutes les cinq minutes en moyenne, dont 30~\% ne se confirment jamais et achètent
2\,000 actions sur le marché visible à la place. Ce qui reste après 30 minutes est acheté immédiatement sur le marché visible. Un sondeur interroge le dark pool toutes les deux secondes avec 100 actions,
en alternant les sens~; lorsque deux sondes consécutives d'un même sens sont exécutées dans les cinq secondes, il achète (ou vend) 500 actions sur le marché visible devant l'ordre qu'il a repéré.

\omcode{code/microstructure/09-dark-pools-and-blocks/python/mx_dark.py}{125}{133}{Le sondeur~: une sonde exécutée est notée, et deux consécutives déclenchent une transaction par anticipation sur le marché visible.}\label{lst:mx:dark-pools-and-blocks:ping}

Les coûts observés sont bruités, si bien que la simulation rend compte aussi des deux coûts qu'elle connaît. Chaque achat visible de $q$ fait bouger $P^\ast$ de $\lambda q$ et relève le
prix de chaque action que le fonds achète ensuite~: sommé sur les ordres visibles enfants du fonds, cela donne son coût d'impact propre, $\lambda Q/2\,(1-1/n)=4.95$ points
de base pour 100 enfants~; sommé sur les transactions du sondeur et des vendeurs qui fuient, qui devancent l'ordre, lorsqu'elles ont été déclenchées par le fonds, cela donne le coût de fuite.
Sur seize graines avec les mêmes événements exogènes pour chaque stratégie (\cref{fig:mx:dark-pools-and-blocks:costs})~:

\begin{itemize}
\item \emph{Visible}~: shortfall 3.8~points de base (erreur type 1.4), dérive du mid sur la fenêtre 7.3 (2.3), impact propre 4.95, aucune fuite.
\item \emph{Dark}~: l'\omterm{def:mx:order-types-and-their-uses:peg}{ordre indexé} se complète en environ 360~secondes, mais le sondeur le repère~: 88 de ses sondes s'exécutent contre le fonds, qui achète 44\,000 actions par anticipation,
et la fuite coûte 10.9~points de base (0.3). Shortfall 12.5 (0.5), dérive 26.1 (2.2).
\item \emph{Dark avec minimum de 1\,000}~: aucune sonde ne peut l'exécuter et il n'y a pas de fuite, mais seulement 46~\% de l'ordre s'exécute dans le dark pool~; le reste est acheté
à la fin, pour 1.4~points de base d'impact propre. Shortfall 9.5 (1.6), dérive 11.9 (1.6).
\item \emph{Bloc}~: 79~\% se croisent en blocs, 0.75 invitations qui fuient par simulation coûtent 0.6~points de base. Shortfall $-2.9$ (1.9), dérive $-1.7$ (2.2).
\end{itemize}

\begin{omfigure}
\begin{tikzpicture}
\begin{axis}[width=10cm,height=5.4cm,ybar,bar width=9pt,xtick={0,1,2,3},xticklabels={visible,dark,{dark, min 1\,000},bloc},
  ylabel={\small points de base},tick label style={font=\footnotesize},ymin=-6,ymax=30,grid=major,grid style={gray!25},
  legend columns=3,legend style={font=\footnotesize,at={(0.5,-0.2)},anchor=north,draw=none,fill=none,
  /tikz/every even column/.append style={column sep=8pt}},table/col sep=comma,enlarge x limits=0.2]
\addplot[omDef,fill=omDef!40,error bars/.cd,y dir=both,y explicit] table[x=k,y=shortfall,y error=shortfall_se]{figdata/microstructure/09-dark-pools-and-blocks/costs.csv};
\addplot[omThm,fill=omThm!40] table[x=k,y=leak]{figdata/microstructure/09-dark-pools-and-blocks/costs.csv};
\addplot[omProp,fill=omProp!40] table[x=k,y=own]{figdata/microstructure/09-dark-pools-and-blocks/costs.csv};
\legend{shortfall,fuite (comptabilisée),impact propre (comptabilisé)}
\end{axis}
\end{tikzpicture}

\omcaption{Un achat de 20\,000 actions exécuté de quatre façons dans le marché simulé~: shortfall par rapport au mid d'arrivée (moyenne et erreur type sur seize graines),
et les deux coûts que la simulation peut comptabiliser exactement. Un ordre laissé au repos sans protection dans le dark pool fuit davantage que ce que coûte la négociation sur le carnet visible. Données~:
\texttt{mx\_dark.compare}.}\label{fig:mx:dark-pools-and-blocks:costs}
\end{omfigure}

Les leçons tiennent au-delà des chiffres du modèle. Un ordre dark au repos est visible de quiconque le sonde, et chaque sonde coûte 100 actions à son émetteur au
mid~; la fuite peut coûter davantage que l'impact que la plateforme dark a été choisie pour éviter. Une \omterm{def:mx:dark-pools-and-blocks:minqty}{quantité minimale d'exécution} supérieure à la taille de la sonde arrête la fuite et coûte
l'exécution complète. Une exécution au mid économise le spread mais pas l'impact lorsque la contrepartie est un flux qui, sinon, aurait vendu sur le carnet visible~: la prendre
hors du marché visible déplace le prix de la même façon. Les blocs sont bon marché dans la simulation parce que leurs vendeurs apportent un volume qui, sinon, n'aurait pas
été négocié~: le coût d'une transaction est le déséquilibre qu'elle crée, où qu'elle soit publiée.

\section{Dark pools dans les dossiers de sanction}

Le secret est le produit, et les dossiers de sanction montrent ce qui se passe lorsque les opérateurs ne le gardent pas. Le 12 août 2015, ITG et sa filiale AlterNet
Securities ont accepté de verser 20.3~millions de dollars, en reconnaissant qu'un desk pour compte propre non divulgué, Project Omega, avait accédé aux flux en direct de son dark pool
aux ordres et aux exécutions de ses abonnés, et y a négocié contre eux dans POSIT. Le 31 janvier 2016, Barclays a accepté de verser 70~millions de dollars, dont la moitié à la SEC et
la moitié au New York Attorney General, pour avoir déformé la manière dont son Liquidity Profiling surveillait son pool LX, et Credit Suisse a payé 84.3~millions de dollars au titre de son
pool Crossfinder, où elle a déformé son Alpha Scoring des flux des abonnés, a accepté plus de 117~millions d'ordres sub-penny et a averti deux
sociétés de haute fréquence des ordres de ses clients via son service Crosslink. Chacun constitue un cas de \omterm{def:mx:dark-pools-and-blocks:leakage}{fuite d'information} imputable à la plateforme elle-même.

\begin{dated}{2026-09}{Plafonds du dark trading dans l'Union européenne}\label{dat:mx:dark-pools-and-blocks:svc}
La révision de MiFIR est entrée en vigueur le 28 mars 2024. Elle remplace le double plafond de volume sur le dark trading par un plafond de volume unique qui ne s'applique
uniquement à la dérogation au prix de référence, à compter de 18~mois après son entrée en vigueur~; l'ESMA prévoyait d'ouvrir son premier dossier de suspension au titre du nouveau plafond pour le 9 octobre 2025.
\end{dated}

\section{Tutoriel~: le bloc qui a fui}\label{tut:mx:dark-pools-and-blocks:tut}

\begin{tutorial}
\textbf{Objectif.} Construire une plateforme visible et une plateforme dark dotée d'une couche de blocs, mesurer qui s'exécute dans le dark, et chiffrer la fuite d'un ordre de grande taille.
\textbf{État final~:} la \cref{fig:mx:dark-pools-and-blocks:costs} et les chiffres des sections 3 et 4.
\begin{tutsteps}
\item \textbf{Plateformes.} \texttt{mx\_dark.Market(seed, strategy)}~: deux moteurs pilotés directement, le moteur dark étant évalué par le contrôle N après chaque changement du
meilleur niveau du carnet visible.
\item \textbf{Sélection.} \texttt{zhu()}~: taux d'exécution dans le dark des preneurs de liquidité informés et non informés, capture et mouvement à 30~secondes des exécutions au repos dans le dark et sur
le carnet visible.
\item \textbf{Fuite.} \texttt{run(seed, strategy)} pour \texttt{lit}, \texttt{dark}, \texttt{dark\_min}, \texttt{block}~; \texttt{compare()} sur
seize graines~; tracer avec \texttt{fig\_dark.py}.
\item \textbf{Couche de blocs.} \texttt{firm\_darkpool.ConditionalBook}~: ajouter des conditionnels, collecter les invitations, confirmer, laisser expirer, et lire les \omterm{def:mx:dark-pools-and-blocks:firmup}{taux de confirmation}.
\end{tutsteps}
\textbf{Pour aller plus loin.} Donner au sondeur des sondes de 50 actions et au fonds un minimum de 200~; laisser la plateforme de blocs exclure les participants dont le \omterm{def:mx:dark-pools-and-blocks:firmup}{taux de confirmation}
tombe sous la moitié, et surveiller la fuite.
\end{tutorial}

\section{Construction~: la plateforme dark}\label{bld:mx:dark-pools-and-blocks:darkpool}

\begin{build}
\bfield{Rôle} Liquidité dark et de blocs pour les algorithmes d'exécution des chapitres 16 à 18 et 28, ainsi que pour les mesures de fuite que produit
l'analyse des coûts de transaction du chapitre~19.
\bfield{Interface} \texttt{ConditionalBook(firm\_up\_ns)}~: \texttt{add}, \texttt{cancel}, \texttt{invitations}, \texttt{firm\_up}, \texttt{expire},
\texttt{firm\_up\_rate}~; \texttt{shortfall\_bps}, \texttt{pre\_completion\_drift\_bps}, \texttt{counterparty\_markout}. Dans \texttt{firm.exchsim}~: contrôle
N dans les moteurs Python, C++20 et Rust, \texttt{ExchangeConfig(reference, reference\_ns)}, \texttt{midpoint\_dark\_pool(reference=\dots)}.
\bfield{Règles} Un conditionnel ne s'exécute que lorsque la plus petite des deux tailles atteint les deux minimums~; un ordre ne figure que dans une invitation à la fois~; le bloc ne se croise que
lorsque les deux parties confirment dans le délai imparti~; une invitation expirée compte contre celui qui n'a pas répondu.
\bfield{Tests d'acceptation} \texttt{code/firm/darkpool/tests/}~: une invitation, une confirmation bilatérale, une expiration et les taux qui en découlent~; shortfall, dérive et
mark-out de contrepartie à la main. \texttt{code/firm/exchsim/tests/}~: un \omterm{def:mx:order-types-and-their-uses:peg}{ordre indexé au mid} avec minimum ignore un petit ordre et croise un grand au
mid du carnet visible~; les jeux de test avec trois messages N restent identiques octet pour octet en C++20 et Rust.
\bfield{Pour aller plus loin} Enchères dark périodiques~; notation des participants au niveau de la plateforme, qui bloque les faibles \omterm{def:mx:dark-pools-and-blocks:firmup}{taux de confirmation}~; la courbe de mark-out de contrepartie du chapitre~19.
\end{build}

\begin{omsources}
\item H.~Zhu, «~Do dark pools harm price discovery?~», \emph{Review of Financial Studies} 27(3), 2014.
\item C.~Comerton-Forde et T.~Putni\c{n}\v{s}, «~Dark trading and price discovery~», \emph{Journal of Financial Economics} 118(1), 2015.
\item U.S. Securities and Exchange Commission, communiqué de presse 2015-164 (ITG, AlterNet Securities), 2015~; communiqué de presse 2016-16 (Barclays, Credit Suisse),
2016.
\item European Securities and Markets Authority, \emph{Double Volume Cap Mechanism}, page web, 2025.
\end{omsources}

\section{Exercices}

\begin{exercise}[$\star$]\label{exo:mx:dark-pools-and-blocks:1}
Le marché visible affiche 20.00 au cours acheteur et 20.02 au cours vendeur. Calculer ce que chaque côté d'un croisement au mid de 10\,000 actions économise par rapport au franchissement du spread sur le carnet visible.
\end{exercise}

\begin{exercise}[$\star$]\label{exo:mx:dark-pools-and-blocks:2}
Un achat au mid au repos de 5\,000 actions, avec une \omterm{def:mx:dark-pools-and-blocks:minqty}{quantité minimale d'exécution} de 1\,000, rencontre une vente de 300 actions, puis une vente de 1\,500 actions. Déterminer les transactions qui ont lieu.
\end{exercise}

\begin{exercise}[$\star$]\label{exo:mx:dark-pools-and-blocks:3}
Un participant a reçu 40 invitations et en a confirmé 28. Calculer son \omterm{def:mx:dark-pools-and-blocks:firmup}{taux de confirmation} et expliquer pourquoi la plateforme s'y intéresse.
\end{exercise}

\begin{exercise}[$\star\star$]\label{exo:mx:dark-pools-and-blocks:4}
Avec $\lambda=0.0005$ ticks par action sur un titre à 100.00, un sondeur achète 500 actions par anticipation lorsque le fonds a acheté 5\,000 sur 20\,000. Calculer le coût de cette seule
transaction pour le fonds, en points de base de l'ensemble de l'ordre.
\end{exercise}

\begin{exercise}[$\star\star$]\label{exo:mx:dark-pools-and-blocks:5}
Dans la simulation, les preneurs de liquidité informés exécutent 6.3~\% de ce qu'ils tentent dans le dark, et les non informés 17.5~\%. Expliquer l'écart à la lumière de l'argument de Zhu.
\end{exercise}

\begin{exercise}[$\star\star$]\label{exo:mx:dark-pools-and-blocks:6}
Montrer que 100 ordres fils égaux, exécutés sur le marché lit, d'un ordre de 20\,000 actions coûtent $\lambda Q/2\,(1-1/n)$ en impact propre, et le calculer.
\end{exercise}

\begin{exercise}[$\star\star\star$]\label{exo:mx:dark-pools-and-blocks:7}
\emph{Programmation.} Relancer \texttt{dark\_min} avec un minimum de 200 au lieu de 1\,000 sur les seize graines. Comparer la fuite, la part exécutée dans le dark,
le temps de complétion et le shortfall avec les deux stratégies dark.
\end{exercise}

\begin{exercise}[$\star\star\star$]\label{exo:mx:dark-pools-and-blocks:8}
\emph{Trouver la faille.} «~Toutes nos exécutions dans le dark se sont faites au midpoint, donc notre gros ordre ne nous a rien coûté.~»
\end{exercise}

\section{Problème~: le bloc qui a fui}

\begin{problem}[{Problème du week-end --- le bloc qui a fui}]\label{pb:mx:dark-pools-and-blocks:1}
Un fonds doit acheter 20\,000 actions en 30 minutes et demande où envoyer l'ordre. Répondre à l'aide du marché simulé.

\medskip\noindent\textbf{Partie I --- Les plateformes.}
\begin{enumerate}
\item Que rapporte un croisement au midpoint, et d'où vient son prix~?
\item De quoi le moteur dark a-t-il besoin de la plateforme lit, et comment le simulateur la lui fournit-il~?
\item Définir une \omterm{def:mx:dark-pools-and-blocks:minqty}{quantité minimale d'exécution} et donner ses deux effets.
\item Définir un \omterm{def:mx:dark-pools-and-blocks:conditional}{ordre conditionnel} et un \omterm{def:mx:dark-pools-and-blocks:firmup}{taux de confirmation}.
\item Pourquoi une invitation non suivie d'effet fait-elle fuir de l'information~?
\end{enumerate}

\medskip\noindent\textbf{Partie II --- Sélection.}
\begin{enumerate}[resume]
\item Énoncer l'argument de Zhu.
\item Donner les taux d'exécution dans le dark des preneurs de liquidité informés et non informés dans la simulation.
\item Comparer la capture et le mouvement à 30 secondes des exécutions passives dans le dark et dans le carnet lit.
\item Qu'ont établi Comerton-Forde et Putnins au sujet des transactions de bloc~?
\end{enumerate}

\medskip\noindent\textbf{Partie III --- Quatre façons d'acheter.}
\begin{enumerate}[resume]
\item Décrire le sondeur et les vendeurs qui laissent fuir l'information.
\item Établir le coût d'impact propre de la stratégie lit.
\item Comment le coût de la fuite est-il comptabilisé~?
\item Donner le shortfall, la dérive et les coûts comptabilisés des quatre stratégies.
\item Pourquoi l'ordre dark non protégé coûte-t-il plus cher que l'ordre lit~?
\item Que corrige la quantité minimale de 1\,000, et que coûte-t-elle~?
\end{enumerate}

\medskip\noindent\textbf{Partie IV --- Plateformes et verdict.}
\begin{enumerate}[resume]
\item Pourquoi les blocs sont-ils bon marché dans la simulation, et quand cesseraient-ils de l'être~?
\item Résumer les cas ITG, Barclays et Credit Suisse en une phrase chacun.
\item Qu'est-ce qui a changé dans le plafond européen sur la négociation dans le dark~?
\item Énoncer le \emph{résultat nommé}~: la dérive de prix avant la complétion d'un ordre parent travaillé dans le dark, face à celle du marché lit, et
le coût de la fuite en points de base.
\item En une phrase~: où le fonds doit-il envoyer son ordre~?
\end{enumerate}
\end{problem}

\section{Questions d'entretien}

\begin{interviewq}[$\star$ \iqroles{trader}]\label{iq:mx:dark-pools-and-blocks:1}
Qu'est-ce qu'un dark pool, et pourquoi une institution en utiliserait-elle un~?
\end{interviewq}

\begin{interviewq}[$\star\star$ \iqroles{trader}]\label{iq:mx:dark-pools-and-blocks:2}
Comment protégeriez-vous un gros ordre passif dans le dark contre toute détection~?
\end{interviewq}

\begin{interviewq}[$\star\star$ \iqroles{chercheur}]\label{iq:mx:dark-pools-and-blocks:3}
Les dark pools nuisent-ils à la \omterm{def:mx:fragmentation-and-routing:discovery}{formation des prix}~? Donnez la théorie et les preuves empiriques.
\end{interviewq}

\begin{interviewq}[$\star\star$ \iqroles{trader, chercheur}]\label{iq:mx:dark-pools-and-blocks:4}
Comment mesureriez-vous la \omterm{def:mx:dark-pools-and-blocks:leakage}{fuite d'information} dans les exécutions de votre société~?
\end{interviewq}

\begin{interviewq}[$\star\star$ \iqroles{développeur}]\label{iq:mx:dark-pools-and-blocks:5}
Concevez la logique d'appariement d'une plateforme de blocs conditionnels. Quel état conserve-t-elle par ordre et par participant~?
\end{interviewq}

\begin{interviewq}[$\star\star\star$ \iqroles{trader}]\label{iq:mx:dark-pools-and-blocks:6}
Vos exécutions au midpoint semblent gratuites, pourtant votre shortfall sur les ordres majoritairement dark est plus élevé que sur les ordres lit. Que se passe-t-il~?
\end{interviewq}
